% Part: normal-modal-logic
% Chapter: filtrations
% Section: introduction

\documentclass[../../../include/open-logic-section]{subfiles}

\begin{document}

\olfileid[zh]{nml}{fil}{int}

\olsection{绪论}

关于一种逻辑，一个重要的问题总是它是否可判定，即是否存在一个能行程序来回答「这个!!{formula}是否有效」这一问题。命题逻辑是可判定的：我们可以通过构造真值表来有效地检验!!a{formula}是否为重言式，而且对于给定的!!{formula}，真值表是有穷的。但我们没有明显的办法检验一个模态公式是否在所有模型中为真，因为模型有无穷多个。我们可以列出与给定公式相关的所有有穷模型，因为只有把世界的子集指派对实际出现在该!!{formula}中的!!{propositional variable}才相关。若可及关系是固定的，则可能的不同指派 $V(p)$ 恰好是~$W$ 的所有子集，而若 $\card{W} = n$，则这样的子集有 $2^n$ 个。若我们的!!{formula}~$!A$ 含有 $m$ 个!!{propositional variable}，那就相应地有 $2^{nm}$ 个不同的、含~$n$ 个世界的模型。对于每一个模型，我们可以通过计算~$!A$ 在每个世界处的真值，来检验~$!A$ 是否在所有世界处为真。当然，我们还得检验所有可能的可及关系，但在~$n$ 个世界上的关系也只有有穷多个（具体地说，其数目就是 $W \times W$ 的子集数，即 $2^{n^2}$。

如果我们感兴趣的并非逻辑~\Log{K}，而是由某种模型类（例如自反传递模型）所定义的逻辑，那么我们还得能够检验可及关系是否属于正确的类型。只要我们所关心的框架能由模态!!{formula}所定义（例如通过检验 \Ax{T} 与 \Ax{4} 是否在框架上有效），我们就能做到这一点。于是，想法就是遍历所有有穷框架，逐一检验它是否属于我们所关心的那类框架，然后列出该框架上的所有可能模型，检验 $!A$ 在每一个模型中是否为真。若不为真，则停止：$!A$ 在所关心的模型类中并非有效。

这个想法有一个问题：我们不知道何时——如果真的有时——可以停止寻找。如果该!!{formula}有一个有穷反模型，我们的程序会找到它。但如果它没有有穷反模型，我们就得不到答案。该!!{formula}可能有效（完全没有反模型），也可能只有一个我们永远不会考察到的无穷反模型。如果我们能证明每一个有反模型的!!{formula}都有有穷反模型，就可以克服这个问题。若情况如此，则称该逻辑具有\emph{有穷模型性}。

但我们如何证明一种逻辑具有有穷模型性呢？一种做法是找到一种把~$!A$ 的无穷（反）模型变为有穷模型的方法。如果能够做到这一点，那么只要存在 $!A$ 在其中不为真的模型，所得的有穷模型也会使 $!A$ 不为真。该有穷模型会出现在我们关于所有有穷模型的清单上，于是对于每一个非有效的公式，我们最终都能确定它不是有效的。若公式有效，我们的程序就不会终止。如果我们还能表明，我们程序提供的有穷模型可以具有某个最大规模，而且该最大规模只取决于!!{formula}~$!A$，那么我们就在寻找反模型的搜索中得到了一个必须检验有穷模型的规模上限。如果到那时还没找到反模型，那就没有反模型。于是，我们的程序实际上就能对任意!!{formula}~$!A$ 判定「$!A$ 是否有效？」这一问题。

把无穷结构变为有穷结构的一种常用策略，是「等同」该结构中在相关方面行为相同的!!{element}。如果在~$\mModel{M}$ 中有无穷多个世界在相关方面行为相同，那么我们或许能把\emph{所有}世界收集到有穷多个（可能是无穷的）这类世界的「类」之中。换言之，我们应该以正确的方式对世界的集合进行划分，即如此划分：使每个划分块包含无穷多个世界，但划分块只有有穷多个。然后定义一个新模型~$\mModel{M^*}$，其世界就是这些划分块。旧模型中的有穷多个划分块给我们带来新模型中的有穷多个世界，即有穷模型。将世界~$w$ 所在的划分块记作 $[w]$。我们希望做到的是 $\mSat{M}{!A}[w]$ 当且仅当 $\mSat{M^*}{!A}[{[w]}]$，因为如果旧模型是~$!A$ 的反模型，我们希望新模型也是~$!A$ 的反模型。这就要求我们以正确方式定义划分，以及~$\mModel{M^*}$ 的可及关系。

要明白这是如何进行的，先设想我们没有可及关系。$\mSat{M}{\Box !B}[w]$ 当且仅当对某个 $v \in W$，$\mSat{M}{\Box !B}[v]$，而对 $\mModel{M^*}$ 亦相同，只是换成 $[w]$ 与 $[v]$。作为最初的构想，称两个世界 $u$ 与 $v$ 是等价的（属于同一划分块），若它们在~$\mModel{M}$ 中对所有!!{propositional variable}都一致，即 $\mSat{M}{p}[u]$ 当且仅当 $\mSat{M}{p}[v]$。令 $V^*(p) = \Setabs{[w]}{\mSat{M}{p}[w]}$。我们的目标是证明 $\mSat{M}{!A}[w]$ 当且仅当 $\mSat{M^*}{!A}[{[w]}]$。显然，我们会通过归纳来证明这一点：基础情形就是 $!A \equiv p$。首先设 $\mSat{M}{p}[w]$。则按定义 $[w] \in V^*$，所以 $\mSat{M^*}{p}[{[w]}]$。现在设 $\mSat{M^*}{p}[{[w]}]$。这意味着 $[w] \in V^*(p)$，即对某个与~$w$ 等价的 $v$，$\mSat{M}{p}[v]$。但「$w$ 与 $v$ 等价」意指「$w$ 与 $v$ 使所有同样的!!{propositional variable}为真」，所以 $\mSat{M}{p}[w]$。现在谈归纳步骤，例如 $!A \ident \lnot !B$。则 $\mSat{M}{\lnot !B}[w]$ 当且仅当 $\mSat/{M}{!B}[w]$ 当且仅当（由归纳假设）$\mSat/{M^*}{!B}[{[w]}]$ 当且仅当 $\mSat{M^*}{\lnot !B}[{[w]}]$。对其余非模态算子亦类似。它对 $\Box$ 也成立：设 $\mSat{M^*}{\Box !B}[{[w]}]$。这意味着对每个 $[u]$，$\mSat{M^*}{!B}[{[u]}]$。由归纳假设，对每个 $u$，$\mSat{M}{!B}[u]$。于是，$\mSat{M}{\Box !B}[w]$。

在一般情形中，我们还必须为 $\mModel{M^*}$ 定义可及关系，事情就更复杂了。若~$\mModel{M^*}$ 的可及关系~$R^*$ 满足使上述归纳证明得以进行所需的条件，则称模型~$\mModel{M^*}$ 为一个\emph{过滤}。于是任意过滤~$\mModel{M^*}$ 都会使 $!A$ 在 $[w]$ 处为真，当且仅当 $\mModel{M}$ 使 $!A$ 在~$w$ 处为真。然而，我们现在还得表明过滤\emph{存在}，即我们能定义~$R^*$，使其满足所需的条件。不过，为了使这一点成立，我们必须要求世界 $u$、$v$ 被视为等价的依据，不仅是它们在所有!!{propositional variable}上一致，还要在~$!A$ 的所有子-!!{formula}上一致。由于 $!A$ 只有有穷多个子-!!{formula}，这仍能保证该过滤是有穷的。定义过滤的方法不止一种，而为了确保过滤的可及关系满足所需的性质（例如自反性、传递性等），我们必须在~$R^*$ 的定义上多用心思。


\end{document}
